\chapter{Building, Testing, and Booting Hydra}
\label{app:building}

This appendix provides a practical, comprehensive engineering guide for compiling the Hydra operating system from source code, assembling bootable UEFI media, and testing the system across modern virtualization hypervisors (Oracle VM VirtualBox and QEMU) and bare-metal hardware.

% ==============================================================================
\section{Toolchain Architecture and Prerequisites}
\label{sec:toolchain_deep}
% ==============================================================================

Hydra is written in modern Rust and targets the bare-metal 64-bit UEFI environment (\texttt{x86\_64-unknown-uefi}). Unlike traditional operating system builds that require complex cross-compiler toolchain bootstraps (such as building custom GCC/binutils cross-targets), Rust natively incorporates the LLVM compiler infrastructure.

\subsection{Installing the Rust Bare-Metal Target}
The host machine requires a stable or nightly Rust toolchain with the UEFI target installed:

\begin{lstlisting}[language=bash, caption={Installing Rust Bare-Metal Targets}]
# Install the x86_64 UEFI target library
rustup target add x86_64-unknown-uefi

# Verify the toolchain installation
rustc --version
cargo --version
\end{lstlisting}

\subsection{Cargo Configuration (\texttt{.cargo/config.toml})}
Hydra configures its build flags to enforce freestanding bare-metal machine code generation:
\begin{lstlisting}[caption={Cargo Configuration for Hydra (\texttt{.cargo/config.toml})}]
[build]
target = "x86_64-unknown-uefi"

[target.x86_64-unknown-uefi]
rustflags = [
    "-C", "link-arg=/ENTRY:efi_main",
    "-C", "link-arg=/SUBSYSTEM:EFI_APPLICATION",
    "-C", "panic=abort",
    "-C", "code-model=large",
]
\end{lstlisting}
\begin{itemize}
    \item \texttt{/SUBSYSTEM:EFI\_APPLICATION}: Generates a standard PE32+ executable header recognized by UEFI firmware.
    \item \texttt{panic=abort}: Eliminates stack-unwinding landing pads, producing compact and deterministic kernel machine code.
    \item \texttt{code-model=large}: Ensures that 64-bit absolute addresses can span the entire canonical higher-half address space.
\end{itemize}

% ==============================================================================
\section{Compiling the Microkernel and Userland Servers}
\label{sec:compiling_deep}
% ==============================================================================

Compiling the entire operating system is orchestrated through the root build automation:

\begin{lstlisting}[language=bash, caption={Building the Operating System from Source}]
# 1. Compile the Ring 0 Microkernel
cargo build --release --target x86_64-unknown-uefi

# 2. Stage the UEFI Boot Binary
mkdir -p dist/x86_64 esp/EFI/BOOT
cp target/x86_64-unknown-uefi/release/hydra.efi dist/x86_64/bootx64.efi
cp target/x86_64-unknown-uefi/release/hydra.efi esp/EFI/BOOT/BOOTX64.EFI
\end{lstlisting}

% ==============================================================================
\section{Packaging the Bootable FAT32 and Hybrid ISO Media}
\label{sec:packaging_deep}
% ==============================================================================

UEFI firmware does not boot from legacy MBR master boot code. Instead, it reads a standard FAT32 partition containing the file \texttt{\textbackslash EFI\textbackslash BOOT\textbackslash BOOTX64.EFI}.

Hydra's root \texttt{Makefile} builds both a raw FAT32 hard disk image and a hybrid El Torito CD/DVD ISO image:

\begin{lstlisting}[language=bash, caption={Building the Bootable Media Pipeline}]
# Generate the complete distribution media
make build
\end{lstlisting}

The build recipe executes the following discrete operations:
\begin{enumerate}
    \item \textbf{Allocating the Virtual Disk:} Creates a 32~MiB image filled with zeros:
    \begin{verbatim}
    dd if=/dev/zero of=dist/x86_64/hydra.img bs=1M count=32
    \end{verbatim}
    \item \textbf{Formatting the EFI System Partition (ESP):} Formats the disk image as FAT32:
    \begin{verbatim}
    mkfs.vfat -F 32 -n "HYDRA_ESP" dist/x86_64/hydra.img
    \end{verbatim}
    \item \textbf{Populating the Executable Payload:} Uses \texttt{mtools} to copy the bootloader and Ring~3 driver ELF binaries without requiring root loopback mounts:
    \begin{verbatim}
    mmd -i dist/x86_64/hydra.img ::/EFI ::/EFI/BOOT ::/bin
    mcopy -i dist/x86_64/hydra.img dist/x86_64/bootx64.efi ::/EFI/BOOT/BOOTX64.EFI
    mcopy -i dist/x86_64/hydra.img esp/bin/* ::/bin/
    \end{verbatim}
    \item \textbf{Constructing the El Torito Hybrid ISO:} Uses \texttt{xorriso} to generate a bootable hybrid ISO:
    \begin{verbatim}
    xorriso -as mkisofs -R -f \
        -e EFI/BOOT/BOOTX64.EFI -no-emul-boot \
        -isohybrid-gpt-basdat \
        -o dist/x86_64/kernel.iso esp/
    \end{verbatim}
\end{enumerate}

% ==============================================================================
\section{Running Hydra in VirtualBox: The Automated Pipeline}
\label{sec:virtualbox_pipeline}
% ==============================================================================

On Windows workstations, Hydra provides the \texttt{vbox.bat} automation harness. The script automates VirtualBox VM configuration, storage controller attachment, and media registry UUID synchronization:

\begin{lstlisting}[language=bash, caption={VirtualBox Execution Commands}]
# Full compile, package, and launch
.\vbox.bat

# Fast launch skipping kernel compilation
.\vbox.bat --no-build

# Resynchronize ZFS pool disk image from template
.\vbox.bat --sync
\end{lstlisting}

\subsection{VirtualBox Hardware Configuration Details}
The \texttt{vbox.bat} script configures the VM using Oracle's \texttt{VBoxManage} command-line interface:
\begin{itemize}
    \item \textbf{Firmware:} 64-bit UEFI firmware enabled (\texttt{VBoxManage modifyvm Hydra --firmware efi}).
    \item \textbf{Processors and Memory:} 2 virtual CPU cores and 4096~MiB (4~GiB) of base RAM.
    \item \textbf{Storage Architecture:} PIIX4 IDE Controller:
    \begin{itemize}
        \item Primary Master (Port 0, Device 0): \texttt{hydrapool.vdi} (16~GiB OpenZFS storage pool disk).
        \item Secondary Master (Port 1, Device 0): \texttt{kernel.iso} (Bootable hybrid CD-ROM image).
    \end{itemize}
    \item \textbf{Network Controller:} Intel PRO/1000 MT Desktop (\texttt{82540EM}) in NAT mode with cable connected and zero link-up delay.
    \item \textbf{Serial Logging:} COM1 mapped directly to a raw host file (\texttt{scratch\textbackslash vbox\_serial.log}), capturing every line of microkernel debug output.
\end{itemize}

% ==============================================================================
\section{Running Hydra in QEMU: The Developer Workflow}
\label{sec:qemu_workflow}
% ==============================================================================

For rapid developer iteration and automated Continuous Integration (CI), Hydra runs under **QEMU**:

\begin{lstlisting}[language=bash, caption={Launching Hydra under QEMU}]
qemu-system-x86_64 \
    -m 4096 \
    -smp 2 \
    -bios /usr/share/ovmf/OVMF.fd \
    -drive file=dist/x86_64/hydra.img,format=raw \
    -drive file=dist/x86_64/hydrapool.qcow2,format=qcow2,if=ide \
    -netdev user,id=net0 -device e1000,netdev=net0 \
    -serial stdio \
    -vga std
\end{lstlisting}

\subsection{Remote GDB Kernel Debugging}
To debug microkernel panics or inspect page tables dynamically, QEMU can be started with the GDB stub enabled:
\begin{lstlisting}[language=bash, caption={Debugging with QEMU and GDB}]
# Start QEMU and freeze CPU at entry point
qemu-system-x86_64 -s -S [options...]

# In another terminal, connect with GDB
gdb target/x86_64-unknown-uefi/release/hydra.efi
(gdb) target remote localhost:1234
(gdb) break src/main.rs:139
(gdb) continue
\end{lstlisting}

This setup allows engineers to single-step machine instructions across the \texttt{CR3} reload boundary, inspect hardware control registers, and examine memory layouts with complete symbolic fidelity.

% ==============================================================================
\section{Advanced GDB Scripting and Page Table Inspection}
\label{sec:gdb_scripting}
% ==============================================================================

When debugging complex virtual memory faults (such as unexpected Page Faults during Copy-on-Write cloning or higher-half address translations), developers can automate GDB using Python and custom macro scripts:

\begin{lstlisting}[caption={Automated GDB Configuration Script (\texttt{.gdbinit})}]
# Set CPU architecture to x86-64
set architecture i386:x86-64
set disassembly-flavor intel

# Connect to QEMU GDB server
target remote localhost:1234

# Define helper to walk a 4-level PML4 page table from GDB
define walk_pml4
    set $vaddr = $arg0
    set $cr3 = $cr3 & ~0xFFF
    set $pml4_idx = ($vaddr >> 39) & 0x1FF
    set $pdpt_idx = ($vaddr >> 30) & 0x1FF
    set $pd_idx   = ($vaddr >> 21) & 0x1FF
    set $pt_idx   = ($vaddr >> 12) & 0x1FF
    printf "Virtual Address: 0x%lx\n", $vaddr
    printf "  PML4 Index: %d, PDPT: %d, PD: %d, PT: %d\n", \
           $pml4_idx, $pdpt_idx, $pd_idx, $pt_idx
end

# Automatically break on microkernel panic handler
break src/console.rs:panic_handler
\end{lstlisting}

% ==============================================================================
\section{VirtualBox Media Registry Management and UUID Resolution}
\label{sec:vbox_uuid_resolution}
% ==============================================================================

A frequent operational hurdle encountered when developing operating systems inside Oracle VM VirtualBox is the **Virtual Medium Registry Conflict**.

\subsection{The Medium Registry Architecture}
VirtualBox maintains an internal global XML database (\texttt{C:\textbackslash Users\textbackslash <user>\textbackslash.VirtualBox\textbackslash VirtualBox.xml}) containing unique UUID signatures for every registered hard disk image (\texttt{.vdi}).

When a build script recreates or repacks the OpenZFS storage pool disk (\texttt{dist/x86\_64/hydrapool.vdi}), the file's internal VDI header generates a new random UUID. If VirtualBox attempts to launch the VM while its XML registry still references the old UUID, the hypervisor aborts startup with the classic COM error:
\begin{lstlisting}[caption={VirtualBox UUID Conflict Error Message}]
Failed to start the virtual machine Hydra.
UUID {bd6bcf20-4d01-4f23-be03-85e89ec2d843} of the medium 
'D:\os\dist\x86_64\hydrapool.vdi' does not match the value 
{c74cb2ef-2258-4a8a-9319-75f0cd4ab159} stored in the media registry 
('C:\Users\<user>\.VirtualBox\VirtualBox.xml').
Result Code: E_FAIL (0x80004005)
Component:   MediumWrap
\end{lstlisting}

\subsection{Automated UUID Repair Protocol}
To solve this conflict automatically without manual XML editing, Hydra's \texttt{vbox.bat} execution script implements an automated unregister-and-reattach pipeline:
\begin{lstlisting}[language=bash, caption={VBoxManage UUID Automated Synchronization}]
@echo off
set VBOXMANAGE="C:\Program Files\Oracle\VirtualBox\VBoxManage.exe"

echo [*] Synchronizing VirtualBox Media Registry UUIDs...

:: 1. Detach disk from VM storage controller
%VBOXMANAGE% storageattach Hydra --storagectl "IDE Controller" ^
    --port 0 --device 0 --type hdd --medium none >nul 2>&1

:: 2. Close medium and purge stale UUID from media registry
%VBOXMANAGE% closemedium disk "D:\os\dist\x86_64\hydrapool.vdi" >nul 2>&1

:: 3. Re-attach disk with fresh UUID signature
%VBOXMANAGE% storageattach Hydra --storagectl "IDE Controller" ^
    --port 0 --device 0 --type hdd --medium "D:\os\dist\x86_64\hydrapool.vdi"

:: 4. Launch VM
%VBOXMANAGE% startvm Hydra
\end{lstlisting}

% ==============================================================================
\section{Headless Continuous Integration and Test Harness}
\label{sec:ci_test_harness}
% ==============================================================================

Hydra features a completely headless test automation pipeline executed via \texttt{make test}. The pipeline launches QEMU with the \texttt{-display none -serial stdio} flags, redirecting microkernel serial logs to a python test runner. The test runner monitors serial output for milestone markers:
\begin{enumerate}
    \item \texttt{[+] Milestone 1: Paging and Intrusive Frame Allocator OK}
    \item \texttt{[+] Milestone 2: Preemptive Multitasking and Context Switch OK}
    \item \texttt{[+] Milestone 3: OpenZFS Storage Pool Imported (Uberblock Ring Verified)}
    \item \texttt{[+] Milestone 4: Isolated Ring 3 Drivers Online (ATA, e1000, Display)}
    \item \texttt{[+] Milestone 5: Linux Compatibility ABI Operational (glibc/musl)}
\end{enumerate}
If all milestones are reached within a 30-second timeout, the test runner signals success (exit code 0) and issues an ACPI poweroff command via the serial port.
